\section{Interchange and sync}\label{sec:interchange}

Two more outward faces remain, both about \emph{transport} rather than trust: how the record is read by
tooling that was not built for it, and how it travels with the code it explains. Both are deliberately
narrow. A provenance projection carries lineage and attribution and nothing else; a git binding carries the
trace's identity and nothing else. What the record \emph{means} --- freshness, contested-ness, the residual
surface, the verdicts --- stays with \code{ponens}, where it can be recomputed.

\subsection{Interchange: projecting to W3C PROV}\label{sec:interchange-prov}

An agent's provenance ought to be legible to tools that were not built for our format. We therefore project
\trace{} onto \textbf{W3C PROV}~\cite{prov}, the standard vocabulary for provenance
(\code{ponens trace export --to prov}, emitting PROV-JSON). The mapping is close to structural: a typed
artifact becomes a \code{prov:Entity}, an action a \code{prov:Activity}, the assistant (with its model) a
\code{prov:Agent}, and the trace itself a \code{prov:Bundle}. The lineage edges carry over one-to-one ---
\code{producer\_action\_id} to \code{wasGeneratedBy}, an action's inputs to \code{used},
\code{derived\_from} to \code{wasDerivedFrom}, and \code{supersedes} to a \code{wasDerivedFrom} of type
\code{prov:Revision} (i.e.\ \code{wasRevisionOf}). Every artifact and action is \code{wasAttributedTo} /
\code{wasAssociatedWith} the agent. Edges whose endpoints are not present in the trace (a dangling
\code{derived\_from}) are dropped rather than emitted, so the result is a well-formed PROV instance.

The essential point --- and the reason this belongs beside \S\ref{sec:context} rather than in place of it ---
is that \emph{PROV is an interchange view, not a reasoning state}. PROV models \emph{what was derived from
what, by whom}; it has no vocabulary for the evidence-logic certificate that makes the trace a context in the
first place. Concretely, the projection is \textbf{faithful for lineage and attribution} and \textbf{lossy for
judgement}:

\begin{itemize}
\item \textbf{Freshness} (I4) is \emph{derived} by \code{ponens} from a dependency-closure fingerprint; PROV's
\code{wasInvalidatedBy} records an invalidation \emph{event someone asserts}, not a \emph{computed} staleness,
so it is not emitted.
\item \textbf{Residuals and defeaters} --- the negative space and the contested-ness of a claim --- have no
PROV relation. A \code{Residual} is exported as a \code{prov:Entity} with its \code{kind}, \code{severity},
and \code{status} riding as \code{ponens:} attributes, but PROV assigns them \emph{no meaning}; the blocking of
a contested claim by counter-evidence stays a \code{ponens}-side computation.
\item \textbf{Verification verdicts} (proved / refuted / unknown) ride as an opaque \code{ponens:verdict}
attribute on the entity.
\item \textbf{Policies} --- the temporal-logic governance axis --- have no analogue; PROV-CONSTRAINTS is a
fixed consistency checker, not a policy language.
\end{itemize}

The rule of thumb is \emph{lineage and attribution round-trip; judgement does not}. Export is deliberately
one-way: a consumer that needs the certificate --- freshness, contested-ness, the residual surface, the
policy verdict --- reads the \code{ponens} trace, and the standard vocabulary is offered so that provenance
tooling can at least consume chain-of-custody and derivation without learning our format. This is the same
distinction the paper draws throughout: a log records \emph{what happened}; a context carries \emph{the
warrant}. PROV is the log projection of a state that is more than a log.

\subsection{Sync and binding}\label{sec:interchange-sync}

For the warrant to be \emph{useful} to a downstream party, it must travel with the code it explains. Three
facts live in three places and no single layer owns all three: the \textbf{local trace file} is canonical for
\emph{content} (immutable, content-hashed); the \textbf{git commit} is canonical for \emph{version} (which
code state the reasoning is about); and a \textbf{hub} is canonical for \emph{decisions} (status, comments,
review items, snapshots, policy runs). The CLI's job is to bind them and keep them from silently diverging ---
in one line, it makes the hub a git remote for reasoning.

\code{ponens bind} is the load-bearing verb. It stamps the trace with \code{repo} / \code{branch} /
\code{commit\_sha} from git, and writes the trace identity \emph{back} onto the commit --- as a
\code{Trace-Id} trailer or, by default, a \code{git} note (\code{refs/notes/ponens}, non-mutating). This makes
the relationship resolvable in both directions, $\code{commit\_sha} \Leftrightarrow \code{trace\_id}$, so the
reasoning is discoverable from the commit (e.g.\ in a pull-request view) and vice versa. The binding is
\textbf{1:1}: a trace is bound to exactly one commit, and iteration produces a \emph{new} commit and a
\emph{successor} trace rather than overwriting the old one --- the same immutability discipline as
\S\ref{sec:interchange-review}, now mirrored one-to-one with git history. \code{push} / \code{pull} move the
trace and the hub's decision state (\code{push} is idempotent and creates a successor on a content change;
\code{pull} is read-only with respect to the working tree), and \code{status} names the one steady state
(\code{synced}) against the divergences (\code{stale}, \code{dirty}, \code{behind}) that each map to a concrete
next command.

The design decision is \textbf{binding-only}: git carries the trace and the link, and \emph{only} those. This
is what lets the warrant travel with the code \emph{without} committing the churny raw trace state or
reinventing review inside git --- git notes do not merge cleanly and committed review files conflict on every
parallel comment (this is the local regime the merge composition of \S\ref{sec:merge} operates within: git
merges the code and the binding; \code{ponens} re-derives the evidence). Human review of the \emph{diff} stays
on the git platform; computable governance and the frozen, signed \emph{snapshot} (the portable audit artifact
with no platform equivalent) stay on the hub; the commit is the shared anchor connecting them. Let the
platform own diff review, let the hub own governance over the reasoning, and connect them through the commit
--- do not reimplement either inside the other.

